% !TeX root = ../papers/nc-p3.tex
\section{交差ブロック可逆な主枝の全点での正則性}
\label{sec:elliptic-evaluation-obstructions}
定理\ref{thm:klein-elliptic-completion}の全薄係数非零の規格化を用い、
$z\ne1$, $z\lambda^2\ne1$ とする。以下の同時文字置換、始点移動、左右反転は
判定を実行するために使う。分類する同型に始点移動や反転を加える必要はない。
\begin{theorem}\label{thm:elliptic-evaluation-obstructions}
五次元主族の閉包にあるこの表示の窓で、二隣接階数6、一文字全階数を仮定する。
標準Koszul左右AS正則となる必要十分条件は、四辺のいずれを左右いずれから
非零点で評価しても階数が3以上となることである。下界を破る表示は、
辺・向き・文字を選び直した後、次の二型で尽くされる。
\begin{enumerate}
\item $A=\left(\begin{smallmatrix}1&\lambda\\c&d\end{smallmatrix}\right)$,
$B=\left(\begin{smallmatrix}\lambda&f\\1&h\end{smallmatrix}\right)$
であり、節\ref{sec:klein-elliptic-completion}の二切断条件を課す。
\item $s,m\ne0$, $(1-s^2)(1-zm^2)\ne0$ とし、
\[
 T=\begin{pmatrix}1&zm\\m&1\end{pmatrix}^{-1}
       \begin{pmatrix}1&s\\s&1\end{pmatrix},\qquad
 D_0^{\mathsf t}=\operatorname{diag}(1,s)T\operatorname{diag}(1,m)^{-1},
 \qquad
 D_1^{\mathsf t}=\begin{pmatrix}0&1\\s&0\end{pmatrix}
                  T\begin{pmatrix}0&1\\m&0\end{pmatrix}^{-1}.
\]
$A=D_0^{-1}Y_\lambda$, $B=(D_1^{-1}\lambda X_\lambda^{\mathsf t})^{\mathsf t}$。
四つの厚い角が可逆なところで二切断条件は恒等的に成り立つ。
\end{enumerate}
いずれの型も指数増大する商を持つ。この排除には五次元主族の閉包への所属を要しない。
\end{theorem}
\begin{proof}
関係空間の最初の厚い二色の列基底は $\binom{I}{D_0},\binom{I}{D_1}$ で、
$D_0=Y_\lambda A^{-1}$, $D_1=\lambda X_\lambda^{\mathsf t}B^{-\mathsf t}$。
両者は可逆である。薄い二色の係数ベクトルは $(1,1,1,z)$ と $(1,1,1,1)$。
階数1関係において薄い係数を $\alpha,\beta$ とすれば、二つの交差ブロックの
行列式 $\alpha^2-\beta^2,z\alpha^2-\beta^2$ は零であり、$\alpha=\beta=0$。
残った対角二ブロックの一方は零でなければならず、厚い各色の両半分への射影が
同型なので関係全体が零である。巡回・反転しても交差ブロックの可逆性は保たれ、
薄い規格化を取り直せる（反転は行と列を入れ替える）。
全位置で純関係がないので、十分性は定理\ref{thm:involution-quartet-component}による。

低位二座標・高位二座標への評価像の射影を考える。評価点の片方の半分が非零なら、
その側の射影は厚い二色だけで全射である。実際、各色の係数射影は同型で、
下位ビットを入れ替える色も使える。片方だけが非零なら薄い関係が反対側の
独立方向を与える。従って階数2以下なら両側が非零で、像は二平面間の同型のグラフである。
色0だけの像をこのグラフの両側に射影すれば、両半分の非零座標数は等しい。
二支持の場合は共通のアフィン文字置換（差1を保つ）で支持を $0,2$ にできる。
全支持の場合も含め $p_0=p_2=1$ としてよい。射影化に加え、四位置の高位の
倍率を同じ数にしテンソル全体の倍率を調整すると、薄い48語を保ったまま
低位と高位の相対倍率を自由に取れるためである。

$p=(1,0,1,0)$ では薄い二行がグラフの行列を $I$ に決める。
両 $D_i$ の第一行は $(1,0)$ でなければならず、これが第一型である。
$p=(1,s,1,m)$ では薄い二行が定理の $T$ の満たす行列方程式を与える。
左の行列が特異なら右の行列も特異で、二行の比例係数が一致するから
$m=s$, $s^2=zm^2=1$ となり $z\ne1$ に反する。厚い二色の評価をこの
グラフに入れると表示した $D_i$ を得る。逆にこれらの式で全評価は階数2となる。

最後に、位置0の双対矢を $kp$、位置1を $R_0(p,-)$ の核 $L$、
位置3を $R_3(-,p)$ の核 $H$、位置2を元の四次元全体に制限して周期的な商を取る。
次の計算は八つのSchur等式（第二型）とこの商の関係階数を有理関数体上で確認する。
各辺の色行列の最初の二列（厚色）または一列（薄色）を使う。厚色の最初の二列は
交差ブロックにより独立である。$E$ を両端それぞれの薄い二行の評価行列とすると
$\binom{-E_{\rm low}^{-1}E_{\rm high}}{I}$ が核の基底であり、厚い四行もこの基底を消す。
第一型で $L=\binom{-I}{I}$, $H=\binom{\operatorname{diag}(-1,-1/\lambda)}{I}$。
第二型で消去の分母は $\lambda,1-zm^2$ の単元倍の積である。
これらの基底を同じ記号で表し、中央の辺の関係行列を $r$ とする。
$L^{\mathsf t}r$ と $rH$ を一列に並べ、それぞれ右・左に二次元の単位基底を
テンソルして得る $16\times24$ 行列の階数は、第一型の独立六変数
$c,d,f,h,z,\lambda$ 上で14、第二型の四変数 $s,m,z,\lambda$ 上で12である。
\path{computations/verify_elliptic_evaluation_obstructions.py} はこの行列を生成して厳密に掃き出す。
第一型では第二切断の条件を課す前から選んだ列についての上界が成立し、
実際の窓に制限したときそれらは関係の全基底である。従って例外値でも階数は
この上界を超えない。位置0に接する二辺の制限関係は零なので、
長さ $4n$ の商空間はこの三文字の商の $n$ 回のテンソル積に一次元因子を掛けたものとなる。
次元は少なくとも $2^n$ または $4^n$ で、標準Hilbert関数に反する。
\end{proof}
